\section{ASK Problem Formulation}
\label{sec:formulation}

We formalize write-time memory governance independently of a controller and
store. The reference protocol supports single-slot versioned, multi-record, and
temporal-event modes. The controlled MGP-6 diagnostic instantiates one atomic
semantic slot so that the write decision can be scored without retrieval or
answer generation.

\subsection{Memory records, state, and types}

A memory record $m$ is a tuple
$m = (\mathit{id}, \tau, v, c, \sigma, \ell, t, p)$, where $\mathit{id}$ is a
unique identifier, $\tau \in \mathcal{T}$ is the memory \emph{type}, $v$ is the
stored value (written in the auditable form \texttt{slot = value}), $c \in [0,1]$
is a stored confidence, $\sigma \in \{\textsf{active}, \textsf{superseded}\}$ is
the record's \emph{status}, $\ell$ is the supersession pointer (\texttt{supersedes\_id},
possibly null), $t$ is a validity timestamp, and $p$ is the record's
\emph{provenance} (the source event/turn that established it). We use three
types, ordered by expected durability:
\[
\mathcal{T} = \{\textsf{stable\_pref},\ \textsf{episodic\_fact},\ \textsf{transient\_pref}\}.
\]
A \emph{memory state} $M$ is a set of records. In single-slot mode a well-formed
state has at most one active record for the slot; an empty initial state is
allowed. Every non-null $\ell$ points to a retained superseded predecessor.
Multi-record mode permits compatible active records. The reference protocol
encodes provenance and evidence digests in its records and events. The evaluated
SQLite runtime stores the core type, status, validity, and lineage fields, while
source session and turn provenance remain in the pipeline trace rather than a
fully normalized provenance table.

\subsection{Candidate update and evidence}

At a write step the agent extracts a \emph{candidate update} $u$ (a proposed new
value with a type guess) from an event against the current state $M$. Let $E$
be the attributable evidence associated with $u$, and let
$m^{\star}=R(M,u,E)$ be the prior record selected by a target-resolution
function $R$, if one exists. The controller can use
extraction confidence, the specificity and recency of the supporting span, the
presence of negation or hedging markers, and the classified logical relation of
$u$ to the active record (support, contradiction, or independence) as evidence
signals. The formulation is agnostic to candidate extraction, target resolution,
and relation-classification methods.

\subsection{Action space and transition function}
\label{sec:actionspace}

The controller $\pi$ emits
$a=\pi(M,u,E,m^{\star}) \in \mathcal{A} =
\{\keep, \supersede, \abstain\}$, which induces a deterministic transition
$M' = \delta(M,u,m^{\star},a)$:
\begin{itemize}[leftmargin=1.5em,topsep=2pt,itemsep=1pt]
\item Under \keep{}, the transition never invalidates the prior active record.
In a single-slot diagnostic, a durable candidate leaves the state unchanged;
a clear temporary override may be added only as a separate
\textsf{transient\_pref}. In a multi-record runtime, a compatible candidate may
instead be appended as another active record while the prior remains active.
Thus \keep{} means ``admit or retain without replacement,'' not necessarily
``perform no insertion.''
\item $\delta(M, u, m^{\star}, \supersede)$ marks the prior active record
\textsf{superseded} and adds one \textsf{active} successor whose $\ell$ points back
at it. The prior value is \emph{retained}, not deleted, and no longer active in
the slot; exactly one active record of the successor's type remains.
\item $\delta(M, u, \abstain) \equiv M$ at the modeled field level: the state is
byte-for-byte equivalent. \abstain{} is a decision to \emph{not commit} the
candidate, distinct from committing a null value.
\end{itemize}
A transition is an \emph{invalid commit} if it changes the active state on an item
whose gold action is \abstain, or otherwise violates a state-transition invariant
(Section~\ref{sec:mgp6}). This is the failure the abstention and provenance
machinery is designed to prevent.

Table~\ref{tab:action-example} illustrates the three actions while holding the
standing value and candidate fixed. The evidence, rather than the candidate
string, determines the warranted transition.

\begin{table}[t]
\centering
\small
\caption{One candidate under three evidence conditions. The standing record is
\texttt{status=available}; the candidate is \texttt{status=do-not-disturb}.}
\label{tab:action-example}
\begin{tabular}{@{}p{0.19\linewidth}p{0.34\linewidth}p{0.37\linewidth}@{}}
\toprule
Action & Evidence reading & State consequence \\
\midrule
\keep{} & A temporary do-not-disturb period expired and
\texttt{available} was reasserted later. & Retain the active standing record;
do not invalidate it for the stale candidate. \\
\supersede{} & A durable do-not-disturb change occurred after the temporary
period ended. & Retain \texttt{available} as superseded, create an active
successor, and link successor to predecessor. \\
\abstain{} & The two timestamps use incomparable time zones and the offsets are
missing. & Leave modeled memory unchanged and record the uncertainty decision
in the audit trace. \\
\bottomrule
\end{tabular}
\end{table}

\subsection{Protocol event and replay contract}

The ASK reference protocol binds each validated transition to an immutable event
$z=(a,\mu,i_p,i_u,i_r,i_{\pi},h_C,h_E,h_M,h_{M'})$, where $\mu$ is the
declared storage mode; $i_p,i_u,i_r,i_{\pi}$ are the compared-prior,
candidate, resulting-record, and predecessor identifiers; $h_C$ and $h_E$ bind
the exact configuration and evidence set; and $h_M,h_{M'}$ bind the before and
after states. An adapter additionally declares whether its trace is sufficient
to replay only controller composition plus transition, or the final policy plus
transition. It must not claim to replay upstream model inference from a logged
final action.

The reference validator requires: \abstain{} makes no commit and preserves
$h_M=h_{M'}$; \keep{} never invalidates the compared prior; \supersede{}
retains that prior, creates an active successor, and resolves its predecessor
edge; identifiers are unique; lineage has neither dangling edges nor cycles;
and a single-slot state has at most one active record. These are reference
protocol correctness conditions, not evidence that the chosen action matches
human semantics. Section~\ref{sec:architecture} distinguishes these normative
events from the narrower decision logs emitted by the evaluated runtime.

\subsection{Four distinctions the formulation makes explicit}
\label{sec:distinctions}

These distinctions are load-bearing and are frequently conflated in memory
systems; we make them explicit because \mgp{} scores each.

\begin{description}[leftmargin=0em,topsep=2pt,itemsep=2pt]
\item[\keep{} vs.\ retrieval absence.] \keep{} is a positive write-time decision
not to invalidate the current active value in the presence of a candidate. It
may append a compatible record in a multi-record store; it is not the same as a
record simply not being retrieved at answer time. \keep{} operates on the
state; retrieval absence operates on a later read.
\item[\abstain{} vs.\ no-op.] A no-op (e.g.\ Mem0's \texttt{NOOP}) fires when the
candidate carries \emph{nothing new}. \abstain{} fires when the candidate carries
something potentially relevant but the evidence is \emph{too weak or ambiguous to
commit} (unreliable source, non-atomic claim, referential uncertainty). Both
leave the state unchanged, but only \abstain{} is an uncertainty-driven refusal;
this is why we report an external \texttt{NOOP}$\to$\abstain{} mapping only as an
upper bound (Section~\ref{sec:results}).
\item[Write-time vs.\ answer-time abstention.] \abstain{} in our formulation is a
\emph{write-layer} non-commit: the store is not modified. It is distinct from an
answer-layer refusal to respond (as scored by end-task benchmarks), which occurs
downstream and does not change memory. \mgp{} withholds the evaluation query
until the write is committed precisely so that answer-time behavior cannot leak
into the write decision.
\item[Supersession vs.\ deletion.] \supersede{} replaces the active value while
\emph{preserving} the prior record as \textsf{superseded} with a lineage edge.
Deletion (and active forgetting) removes information and is explicitly out of
scope for this work; the reserved \textsf{forgotten} status never occurs in our
data.
\end{description}

\subsection{Objective}

Let $y$ be the target action and let $C(a,y)$ be an application-specific cost.
The policy objective is
\begin{equation}
\pi^* = \arg\min_\pi\;\mathbb{E}[C(\pi(M,u),y)],
\end{equation}
subject to state and lineage invariants. The cost is intentionally not assumed
symmetric: committing on ambiguous evidence, missing a valid update, and
unnecessarily abstaining can have different consequences. We therefore report
the full action confusion and expected loss under multiple pre-declared cost
regimes rather than treating accuracy as a universal utility. State validity,
lineage integrity, stale retrieval, and calibration remain separate outcomes;
they are not folded into a post-hoc scalar chosen after seeing test results.
